\documentclass[12pt]{article}
\usepackage{amsmath,amsfonts,fullpage,amsthm,graphicx}

\newtheorem*{theorem}{Theorem}
\newtheorem*{fact}{Fact}
\newtheorem*{goal}{Goal}
\newtheorem*{idea}{Idea}
\newtheorem*{defn}{Definition}
\newtheorem*{ex}{Example}

%Style section
%\setlength{\textheight}{10in}
%\setlength{\textwidth}{7in}
%\hoffset=-0.75in
 \pagestyle{plain}
% \setlength{\topmargin}{0in}
% \setlength{\headsep}{0in}
% \setlength{\oddsidemargin}{0in}
% \setlength{\evensidemargin}{0in}

\begin{document}


\pagestyle{empty} %use this to get rid of page numbers

\noindent
{Math 166 
%- Dr.\ Jessica Striker
 \hskip 1.5 truein  Names:\ \hrulefill}

\noindent
%\vskip 0.3truein
%\bigskip
%\noindent
{A study of series}

\noindent  \hskip 3.5 truein  Section Number:\ \hrulefill
\medskip

\begin{goal} \rm
Find the sum of a geometric or telescoping series, and apply the $n$th term test.
%Find the sum of a geometric or telescoping series, or show the series diverges.
\end{goal}
%\smallskip

\begin{enumerate}
\setcounter{enumi}{-1}
%Name all the tests you can think of for convergence/divergence of series in our series toolbox. Put a * by the tests that only work for series with positive terms.
\item Fill in the following blanks.
\begin{itemize}
\item A \textbf{sequence} is a~ \rule{.8in}{.4pt} ~of numbers.  A sequence converges if its~ \rule{.8in}{.4pt} ~have a limit.
\item A \textbf{series} is a~ \rule{.6in}{.4pt} ~of numbers. A series converges if its~ \rule{.8in}{.4pt} \rule{.2in}{.4pt} \rule{.8in}{.4pt} \rule{.5in}{.4pt} ~has a limit.
\item The {\bf nth term test} tells us that the series $\displaystyle \sum_{n = 1}^\infty a_n$ \rule{1in}{.4pt} if $\displaystyle \lim_{n = 0} a_n \neq$ \rule{.3in}{.4pt}, and if $\displaystyle \lim_{n = 0} a_n = $ \rule{.3in}{.4pt}, it tells us \rule{3in}{.4pt}.
\end{itemize}


\item A \emph{geometric series} is a series of the form, $\displaystyle \sum_{n = 0}^\infty ar^n$ for constants $a$ and $r$.  When does such a series converge?  If the series does converge, what does it converge to?
\vfill

\item Find the fractional expansion of the repeating decimal $0.42424242\overline{42}$ by rewriting it as a geometric series.

\vfill

\newpage

\item Determine whether the series converges or diverges. If it converges, find the sum.

\[1-\frac{1}{\sqrt{3}}+\frac{1}{3}-\frac{1}{3\sqrt{3}}+\frac{1}{9}-\frac{1}{9\sqrt{3}}+\cdots\]
\vfill
\item Determine whether the series $\displaystyle \sum_{n = 1}^\infty \frac{4}{5n}$ converges or diverges.  If it converges, find the sum.
\vfill

\item Determine whether the series $\displaystyle \sum_{n = 5}^\infty \left(1 + \frac{1}{n}\right)^n$ converges or diverges.  If it converges, find the sum.
\vfill
\newpage

\item Determine whether the series $\displaystyle\sum_{n=1}^{\infty}\displaystyle\frac{3}{n(n+2)}$ converges or diverges. If it converges, find the sum. (Hint: Use partial fractions.)

\vfill

\item Determine whether the series $\displaystyle\sum_{n=1}^{\infty}\ln(n) - \ln(n + 1)$ converges or diverges. If it converges, find the sum.

\vfill

\item Determine whether the series $\displaystyle\sum_{n=0}^{\infty}\frac{7 + e^n}{\pi^n}$ converges or diverges. If it converges, find the sum. (Hint: Use limit laws to split into two series.)

\vfill
%\newpage
%
%\item State the {\bf integral test} in your own words. What three conditions do you need to check in order to use the integral test?
%
%%\vspace{.5in}
%%\item In your team, each person should attempt to use \textbf{one} of the following tests to determine whether the series  converges or diverges: $n$th term divergence test, Comparison test, Limit comparison test, Integral test. Which test(s) worked well for this series? 
%
%\vspace{.7in}
%\item Use either the $n$th term divergence test or the integral test to determine whether the following series converge or diverge. Be sure to verify the hypotheses, state your conclusion, and say which test you used. 
%%You should approach each problem together as a team, but you may each try different tests to more quickly determine what test will work in each problem.
%\begin{enumerate}
%\item $\displaystyle\sum_{n=1}^{\infty} (1.01)^n$
%
%%\vspace{3.5in}
%%\pagebreak
%\vspace{2.5in}
%
%%\item $\displaystyle\sum_{n=1}^{\infty} e^{-3n}$
%%
%%%\vspace{1.5in}
%%%\item $\displaystyle\sum_{n=2}^{\infty} \frac{1}{\sqrt{n^4+n^3}}$
%%
%%\vspace{2.5in}
%%\item $\displaystyle\sum_{n=1}^{\infty} \frac{7^n+1}{2^n}$
%%
%%\vspace{2.5in}
%\item $\displaystyle\sum_{n=2}^{\infty} \frac{n}{n^2+1}$
%
%\vspace{2.5in}
%\item $\displaystyle\sum_{n=1}^{\infty} \frac{\cos(n)}{n^2}$ (Hint: Can you use the alternating series test? Why not? Instead, check for absolute convergence.)
%\item $\displaystyle\sum_{n=1}^{\infty} \frac{1}{4+5^n}$
%\item $\displaystyle\sum_{n=3}^{\infty} \ln(n)$


%\end{enumerate}

\end{enumerate}

\end{document}
